% Capítulo 3: Especificación del Sistema

\section{Definición de requisitos}

\subsection{Requisitos funcionales}

\paragraph{RF-01} El sistema debe permitir seleccionar un punto de corte en la red base (capa
\texttt{initial\_layer}) a partir del cual se aplica el algoritmo Qsimov.

\paragraph{RF-02} El sistema debe enumerar todos los caminos activos entre la capa
\texttt{initial\_layer} y la salida de la red, para capas de tipo Dense, Conv1D/2D/3D y
MaxPooling1D/2D/3D con activación ReLU.

\paragraph{RF-03} Dado un conjunto de muestras de entrenamiento, el sistema debe generar el sistema
de ecuaciones lineales $Ax = b$ donde cada fila representa la contribución de una muestra a los
pesos de los caminos.

\paragraph{RF-04} El sistema debe resolver el sistema lineal mediante back-substitution (solución
exacta) o mínimos cuadrados (solución aproximada).

\paragraph{RF-05} El sistema debe permitir acumular ecuaciones de distintas rondas de entrenamiento
sin resetear el sistema, implementando así el mecanismo de no-olvido.

\paragraph{RF-06} El sistema debe implementar una variante basada en descenso de gradiente sobre
una capa densa restringida a los caminos activos.

\paragraph{RF-07} El sistema debe ser compatible con los frameworks TensorFlow/Keras y PyTorch.

\paragraph{RF-08} El sistema debe permitir guardar y cargar modelos entrenados en un formato propio
(directorios \texttt{.qsi}).

\subsection{Requisitos no funcionales}

\paragraph{RNF-01 (Rendimiento)} Las operaciones de extracción de coeficientes y generación de
caminos deben estar optimizadas con extensiones Cython compiladas.

\paragraph{RNF-02 (Compatibilidad)} La librería debe funcionar con Python 3.8 y las versiones
TensorFlow 2.11 y PyTorch 2.0.

\paragraph{RNF-03 (Modularidad)} Las implementaciones para Keras y PyTorch deben compartir clases
base abstractas que definan la interfaz pública.

\paragraph{RNF-04 (Testeabilidad)} Todas las clases principales deben tener cobertura de tests
unitarios e integración ejecutables con \texttt{pytest}.

\section{Especificación del sistema}

El sistema se organiza en torno a tres componentes principales, cuya relación se describe a
continuación:

\begin{figure}[htbp]
  \centering
\begin{verbatim}
              +-------------------------------+
              |   Red base pre-entrenada      |
              | (Keras Sequential / PyTorch)  |
              +---------------+---------------+
                              |
              +---------------v---------------+
              |         PathSelector           |
              |  - left_model_  (hasta i_layer)|
              |  - right_model_ (desde i_layer)|
              |  - all_paths_                  |
              |  - samples_to_coefficients()   |
              +---------------+---------------+
                              |
          +-------------------+-------------------+
          |                                       |
+---------v-----------+               +-----------v-----------+
|  QsimovLinearSystem |               |    QsimovGradient     |
|  - fit() (acumul.)  |               |  - fit() (por epochs) |
|  - predict()        |               |  - predict()          |
|  - reset_equations()|               |  - model_ (CustomDens)|
+---------------------+               +-----------------------+
\end{verbatim}
  \caption{Diagrama de componentes del sistema Qsimov.}
  \label{fig:componentes}
\end{figure}

\section{Planificación y estimación de costes}

\subsection{Planificación}

El desarrollo del proyecto se ha estructurado en cinco fases principales:

\begin{table}[htbp]
  \centering
  \caption{Planificación del proyecto.}
  \label{tab:planificacion}
  \begin{tabular}{cp{8.5cm}c}
    \toprule
    \textbf{Fase} & \textbf{Descripción} & \textbf{Duración estimada} \\
    \midrule
    F1 & Revisión bibliográfica y estudio del estado del arte                       & 3 semanas \\
    F2 & Comprensión de la librería Qsimov y configuración del entorno experimental & 3 semanas \\
    F3 & Diseño de los experimentos y selección de configuraciones                  & 2 semanas \\
    F4 & Ejecución de experimentos y análisis de resultados                         & 8 semanas \\
    F5 & Redacción de la memoria y preparación de la defensa                        & 3 semanas \\
    \midrule
    \textbf{Total} & & \textbf{19 semanas} \\
    \bottomrule
  \end{tabular}
\end{table}

\subsection{Estimación de costes}

\textbf{Coste de personal} (estimación a precio de junior en prácticas, 15\,€/h; tarifa orientativa
para el sector TIC en España, acorde al convenio colectivo de empresas de ingeniería y consultoría
para titulados recientes):
\begin{itemize}
  \item 19 semanas $\times$ 25 h/semana = 475\,h
  \item Coste: 475\,h $\times$ 15\,€/h = \textbf{7.125\,€}
\end{itemize}

\textbf{Coste de hardware}:
\begin{itemize}
  \item Uso de GPU (NVIDIA, estimación de servidor cloud): $\sim$80\,h $\times$ 0.5\,€/h =
        \textbf{40\,€}
\end{itemize}

\textbf{Software}: Todas las herramientas utilizadas (Python, TensorFlow, PyTorch, Cython) son de
código abierto y gratuitas. Coste: \textbf{0\,€}

\textbf{Coste total estimado: 7.165\,€}
